\section{Zeroing with extended reasoning}
\label{sec:thinking-zeroing}
We transferred seven previously selected interventions, together with an
unmodified baseline and its repeat, to Qwen3-30B-A3B-Thinking-2507. Each arm
was evaluated on MATH500 and GPQA Diamond with thinking enabled and a maximum
of 32,768 generated tokens shared by reasoning and the final answer. The
interventions comprise input embedding weight zeroing at 70\% and 90\%,
input embedding plus output-head weight zeroing at 50\%, MLP intermediate
activation zeroing at 50\% and 80\%, block residual-stream activation zeroing
at 25\%, and MoE output activation zeroing at 40\%. Weight selection is
independent within each token vector, while activation selection is per token
at each selected site. These are inference-time interventions without fine-tuning.
The submitted source registry defines the complete protocol and intervention
configuration; this series uses its own baseline rather than the earlier
instruct checkpoint's baseline.

All 18 evaluations completed, covering 500 mathematics questions and 198
scientific reasoning questions per arm. We checked saved configurations,
frozen dataset manifests, complete question coverage, identical prompts across
arms, and metrics recomputed from saved samples. Final-answer extraction and
correctness were independently rechecked for all 6,282 responses using the
submitted grading code. Grading uses only the final response following the
closing reasoning delimiter. Responses with incomplete reasoning or reaching
the generation cap count as incorrect.

\begin{table}[ht]
\centering
\caption{Accuracy (\%) with the thinking checkpoint and extended generation
budget. MLP, block and MoE output rows zero activations; embedding rows zero
weights.}
\label{tab:thinking-accuracy}
% Transcribed from validated thinking32k benchmarks-summary.json.
\begin{tabular}{lrr}
\toprule
Experiment & MATH500 & GPQA Diamond \\
\midrule
Baseline & 98.60 & 67.68 \\
Baseline repeat & 97.80 & 70.71 \\
Embedding 70\% & 97.80 & 67.68 \\
Embedding 90\% & 98.20 & 70.71 \\
Embedding + output 50\% & 97.80 & 65.66 \\
MLP 0.5 & 98.40 & 68.69 \\
MLP 0.8 & 97.20 & 66.16 \\
Block 0.25 & 97.80 & 65.15 \\
MoE output 0.4 & 98.00 & 70.20 \\
\bottomrule
\end{tabular}

\end{table}

Table~\ref{tab:thinking-accuracy} shows that all selected interventions retain
high MATH500 accuracy: 97.20--98.40\%, compared with 98.60\% for baseline
and 97.80\% for its repeat. Input embedding zeroing at 90\% scores 98.20\%
on MATH500 and 70.71\% on GPQA. MLP activation zeroing at 50\% scores
98.40\% and 68.69\%, respectively. No intervention exceeds the first baseline
on MATH500. On GPQA, the baseline repeat improves from 67.68\% to 70.71\%,
a difference of six questions or 3.03 percentage points. The apparent gains
from embedding 90\% and MoE output 40\% fall within this observed variation;
these single intervention runs do not establish a reliable improvement.

\begin{table}[ht]
\centering
\caption{Mean input and output tokens per question for the thinking series.
Input is the initial prompt including the chat template; output is server-reported
completion usage, including reasoning, the final answer and generated special
tokens. All responses, including incorrect and capped responses, enter the mean.}
\label{tab:thinking-lengths}
% Transcribed from thinking-token-lengths-provenance.json; means rounded to one decimal.
\begin{tabular}{lrrrr}
\toprule
& \multicolumn{2}{c}{MATH500} & \multicolumn{2}{c}{GPQA Diamond} \\
\cmidrule(lr){2-3}\cmidrule(lr){4-5}
Experiment & Input & Output & Input & Output \\
\midrule
Baseline & 102.9 & 5553.1 & 241.8 & 7229.8 \\
Baseline repeat & 102.9 & 5451.7 & 241.8 & 6899.0 \\
Embedding 70\% & 102.9 & 5453.5 & 241.8 & 7245.7 \\
Embedding 90\% & 102.9 & 5345.9 & 241.8 & 7994.7 \\
Embedding + output 50\% & 102.9 & 6692.1 & 241.8 & 10342.7 \\
MLP 0.5 & 102.9 & 5348.3 & 241.8 & 7564.1 \\
MLP 0.8 & 102.9 & 4593.8 & 241.8 & 7086.7 \\
Block 0.25 & 102.9 & 5314.2 & 241.8 & 6711.0 \\
MoE output 0.4 & 102.9 & 5191.9 & 241.8 & 7178.2 \\
\bottomrule
\end{tabular}

\end{table}

Table~\ref{tab:thinking-lengths} separates the two benchmarks rather than
pooling their different question counts. Input length is computed offline using
the pinned checkpoint tokenizer with no additional special tokens. Unlike the
retokenized output text lengths in Table~\ref{tab:output-lengths}, the output
counts here use saved completion usage. They include the full reasoning trace,
not only the final answer. None of these runs uses a forced-answer continuation.
These counts characterize generation workload and do not measure latency or
establish a serving speedup.

Joint input embedding and output-head zeroing at 50\% scores 97.80\% on
MATH500 and 65.66\% on GPQA. It also exhausts the generation budget on 5/500
mathematics responses and 19/198 GPQA responses; respectively five and eighteen
lack a closing reasoning delimiter. Neither GPQA baseline run reaches the cap.
Thus similar aggregate accuracy can coexist with increased generation workload
and more incomplete responses. The capped runs remain in both the accuracy and
length summaries; their observed lengths do not estimate uncapped output length.
Mean output usage for this joint weight intervention increases from 5,553.1 to
6,692.1 tokens on MATH500 (20.5\%) and from 7,229.8 to 10,342.7 tokens on
GPQA (43.1\%) relative to the first baseline. In contrast, MLP activation
zeroing at 80\% reduces MATH500 output usage to 4,593.8 tokens (17.3\%
below baseline), alongside an accuracy decrease of 1.40 percentage points.
Identical prompts give the same mean input length across all nine arms:
102.9 tokens for MATH500 and 241.8 for GPQA.
Comparisons with earlier instruct results change both the checkpoint and budget,
so they cannot isolate the effect of enabling thinking.

\paragraph{Result provenance.}
{\raggedright
The submitted source snapshot is
\texttt{a13fc09b0db6514984e495ab489919f4529ae42d}, whose
\texttt{run.py:experiment\_configs} defines this series. Validated scores and
complete job identifiers are saved in
\path{run/thinking32k/benchmarks-summary.json} and
\path{run/thinking32k/jobs.json}. The checked-in
\path{paper/tables/thinking-token-lengths-provenance.json} records sample and
configuration checksums, the pinned tokenizer and its checksums, and unrounded
length means. Tables transcribe these validated values; accuracies are rounded
to two decimals and lengths to one decimal.
\par}
